\documentclass[a4paper,landscape,10pt]{scrartcl}

\usepackage[T1]{fontenc}
\usepackage{lmodern}
\usepackage[ngerman]{babel}
\usepackage{csquotes}
\usepackage[landscape,margin=10mm,top=9mm,bottom=9mm]{geometry}
\usepackage{xcolor}
\usepackage{multicol}
\usepackage{tabularx}
\usepackage{enumitem}
\usepackage{hyperref}

\pagestyle{empty}
\setlength{\parindent}{0pt}
\setlength{\columnsep}{7mm}
\setlength{\columnseprule}{0.3pt}
\renewcommand{\arraystretch}{1.15}

\definecolor{kopf}{gray}{0.25}
\definecolor{zeile}{gray}{0.94}

\newcommand{\cs}[1]{\texttt{\textbackslash #1}}
\newcommand{\m}[1]{\ensuremath{\langle}\textit{#1}\ensuremath{\rangle}}

\newcommand{\bereich}[1]{%
  \par\addvspace{5pt}%
  {\colorbox{kopf}{\parbox{\dimexpr\linewidth-2\fboxsep}{\color{white}\bfseries\strut #1}}}%
  \par\nopagebreak\vspace{2pt}}

% Zweispaltige Tabelle: Befehl | Wirkung
\newenvironment{befehle}
  {\tabularx{\linewidth}{@{}>{\raggedright\arraybackslash\ttfamily}p{0.47\linewidth}>{\raggedright\arraybackslash}X@{}}}
  {\endtabularx\par}

% langer Befehl: eigene Zeile, Erklärung darunter
\newcommand{\lang}[1]{\multicolumn{2}{@{}l@{}}{\ttfamily #1}\\}

\newcommand{\bsp}[1]{\par{\footnotesize\colorbox{zeile}{\parbox{\dimexpr\linewidth-2\fboxsep}{\ttfamily\raggedright #1}}}\par\vspace{2pt}}

\begin{document}

{\Large\bfseries schooltex -- Cheatsheet}\hfill
{\small \texttt{\textbackslash usepackage\{schooltex\}} \quad·\quad KOMA-Klasse, \texttt{babel}, \texttt{hyperref}, \texttt{enumitem}, \texttt{siunitx} vorher laden \quad·\quad ausführlich: \texttt{schooltex-doku.pdf}}
\par\vspace{3pt}\hrule\vspace{4pt}

\raggedcolumns
\begin{multicols}{3}

% ------------------------------------------------------------
\bereich{Gliederung im Skript}
\begin{befehle}
\textbackslash kap\{\m{Titel}\}   & Kapitel \enquote{1 \textsc{Titel}} \\
\textbackslash uber\{\m{Titel}\}  & Unterkapitel \enquote{1.1 Titel} \\
\textbackslash uuber\{\m{Titel}\} & Unter-Unterkapitel \enquote{1.1.1} \\
\textit{\textrm{jeweils}} *       & ohne Nummer, z.\,B. \cs{uber*}\{\dots\} \\
\end{befehle}

\bereich{Tafelbild}
\begin{befehle}
\textbackslash tafel\{\m{Inhalt}\} & grüne Box \enquote{Tafelbild}, umbrechbar, Absätze erlaubt \\
\end{befehle}

% ------------------------------------------------------------
\bereich{Versuche}
\begin{befehle}
\lang{\textbackslash begin\{experiments\}[\m{Titel}]}
 & Versuchsliste, Titel Standard \enquote{Versuche} \\
\textbackslash V              & Versuch \enquote{V 3.1:} ohne Name \\
\textbackslash V[\m{Name}]    & \enquote{V 3.1 (Name):}, Label \texttt{V:Name}, Eintrag im Verzeichnis \\
\textbackslash Vtext\{\m{Text}\}   & Absatz ohne Nummer in der Liste \\
\textbackslash Vgroup\{\m{Titel}\} & blaue Zwischenüberschrift mit Linie \\
\textbackslash listoftoc\{loe\}    & Versuchsverzeichnis ausgeben \\
\end{befehle}
\bsp{\textbackslash begin\{experiments\}\\
\ \ \textbackslash V[Wellenwanne eben] Ebene Wellen \dots\\
\textbackslash end\{experiments\}}

\bereich{Verweise auf Versuche}
\begin{befehle}
\textbackslash Vref\{A\}         & V 3.2 \\
\textbackslash Vref\{A, B\}      & V 3.2 und 3.4 \\
\textbackslash Vref\{A, B, C\}   & V 3.2, 3.4 und 3.5 \\
\textbackslash Vref*\{A\}        & Versuch 3.2 \\
\textbackslash Vref*\{A, B\}     & Versuche 3.2 und 3.4 \\
\textbackslash Vrefrange\{A\}\{B\}  & V 3.2 bis 3.5 \\
\textbackslash Vrefrange*\{A\}\{B\} & Versuche 3.2 bis 3.5 \\
\end{befehle}
{\footnotesize A, B = Versuchsnamen \emph{ohne} \texttt{V:} und ohne Komma. Ein \texttt{V:} davor schadet nicht.
\cs{ref}\{V:A\} gibt nur die Nummer aus, \cs{autoref}\{V:A\} ergibt \enquote{Versuch 3.2}.}

\bereich{Bilder}
\begin{befehle}
\textbackslash begin\{Vbilder\}[0.45] & Bilder nebeneinander, Breite als Anteil der Zeile \\
\lang{\textbackslash Vbild[\m{Optionen}]\{\m{Datei}\}}
 & Optionen: \texttt{caption=}, \texttt{label=}, \texttt{width=} \\
\lang{\textbackslash minicaption[\m{Label}]\{\m{Text}\}}
 & Teilbild-Beschriftung in einer \texttt{minipage} \\
\lang{\textbackslash minimaincaption\{\m{Text}\}}
 & schließt die Teilbild-Gruppe ab \\
\end{befehle}

% ------------------------------------------------------------
\bereich{Stundenverlauf}
\begin{befehle}
\lang{\textbackslash begin\{plan\}[\m{Titel}]}
 & Tabelle mit eigenem \cs{section}, Standardtitel \enquote{Stundenverlauf} \\
\textbackslash umbruch\{a \textbackslash\textbackslash\ b\} & Zeilenumbruch in einer Zelle \\
\end{befehle}
\bsp{\textbackslash begin\{plan\}\\
\ \ 10' \& Thema \& Ausführung \& \textbackslash lsg \& Material \textbackslash\textbackslash\\
\textbackslash end\{plan\}}
\begin{befehle}
\textbackslash lsg \ \textbackslash lv \ \textbackslash ea & L-S G. \ LV \ EA \\
\textbackslash pa \ \textbackslash ga & PA \ GA \\
\end{befehle}

\bereich{Arbeitsblätter ins Skript einbinden}
\begin{befehle}
\textbackslash declarePath\{\m{Pfad}\} & in der Präambel, Ordner der Arbeitsblätter \\
\textbackslash AB\{\m{Name}\} & bindet \m{Pfad}/\m{Name}/\m{Name}.pdf ein, Label \texttt{pdf:Name} \\
\textbackslash AB[small] & 2 Seiten nebeneinander, quer, mit Rahmen \\
\lang{\textbackslash AB[smalllandscape]}
 & 2 Querseiten untereinander \\
\textbackslash AB[landscape] & jede Seite quer \\
\end{befehle}

% ------------------------------------------------------------
\bereich{Arbeitsblatt / Klassenarbeit}
\begin{befehle}
\lang{\textbackslash newcommand\{\textbackslash Datum\}\{\dots\}}
 & muss vor \cs{header} definiert sein \\
\lang{\textbackslash header\{\m{Min}\}\{\m{Hilfsmittel}\}}
 & Kopf: Bearbeitungszeit, Hilfsmittel, Datum \\
\lang{\textbackslash begin\{Aufgabe\}[\m{Punkte}]}
 & \enquote{Aufgabe n}, Punkte am Rand, werden summiert \\
\textbackslash teilaufgabe\{\m{Text}\} & a), b), c), \dots \\
\textbackslash solution\{\m{Lösung}\} & in der Aufgabe: Lösung speichern \\
\textbackslash printsols & alle Lösungen ausgeben \\
\textbackslash printsols[2] & Lösung zu Aufgabe 2 \\
\textbackslash printsols[last] & Lösung der letzten Aufgabe \\
\textbackslash grid & Karoraster bis Seitenende \\
\textbackslash grid[\m{cm}] & Karoraster mit Höhe in cm \\
\textbackslash Ende & \enquote{Viel Erfolg!} + Punkte/Note \\
\textbackslash Punkte & Gesamtpunktzahl bisher \\
\end{befehle}
\bsp{\textbackslash begin\{Aufgabe\}[4]\\
\ \ \textbackslash teilaufgabe\{Berechne \dots\}\\
\ \ \textbackslash teilaufgabe\{Begründe \dots\}\\
\ \ \textbackslash solution\{\textbackslash teilaufgabe\{\dots\} \textbackslash teilaufgabe\{\dots\}\}\\
\textbackslash end\{Aufgabe\}\\
\textbackslash grid[6]\ \dots\ \textbackslash Ende \ \textbackslash newpage \ \textbackslash printsols}

\end{multicols}
\end{document}
